\label{sec:background}

\subsection{LODES Origin-Destination Data}

The Longitudinal Employer-Household Dynamics (LEHD) program at the U.S.\
Census Bureau compiles the LODES (LEHD Origin-Destination Employment Statistics)
dataset from state Unemployment Insurance wage records matched to employer
and residential address geocodes \citep{abowd2009lodes}.

The LODES dataset has three file types:
\begin{itemize}
  \item \textbf{WAC} (Workplace Area Characteristics): jobs by census block
    of work, used in M.1 \citep{m1econchar}.
  \item \textbf{RAC} (Residence Area Characteristics): workers by census block
    of residence.
  \item \textbf{OD} (Origin-Destination): job counts by home block and work
    block pair.
\end{itemize}
M.4 uses the OD file, which is the most granular LODES product.

\paragraph{OD file structure.}
The OD main file (job type JT00, all jobs) has the following primary columns:
\begin{itemize}
  \item \texttt{h\_geocode} (15-character block GEOID of home location)
  \item \texttt{w\_geocode} (15-character block GEOID of work location)
  \item \texttt{S000} (total number of jobs with this home-work pair)
\end{itemize}
Additional columns (\texttt{SA01}--\texttt{SA03}) break the job count by
age group; these are not used in M.4.
Each row records how many workers live in block \texttt{h\_geocode} and
work in block \texttt{w\_geocode}.
The 2020 vintage is used.
File naming convention:
\begin{center}
  \texttt{\{state\}\_od\_main\_JT00\_2020.csv.gz}
\end{center}
available at
\url{https://lehd.ces.census.gov/data/lodes/LODES8/{state}/od/}.

\paragraph{Block-to-tract aggregation.}
Both \texttt{h\_geocode} and \texttt{w\_geocode} are 15-character block GEOIDs.
The first 11 characters identify the census tract (state 2 + county 3 +
tract 6).
Block-to-tract aggregation proceeds in two steps:
\begin{enumerate}
  \item Truncate \texttt{h\_geocode} and \texttt{w\_geocode} to 11 characters
    to obtain home tract $h$ and work tract $w$.
  \item Sum \texttt{S000} for all block pairs within the same $(h, w)$ tract pair.
\end{enumerate}
The result is a sparse matrix $F$ where $F[h, w]$ is the total number of
workers residing in home tract $h$ and employed in work tract $w$.

\paragraph{File size and sparsity.}
The OD file for a large state can be substantial.
The 2020 NC OD file has approximately 2.1 million rows before tract
aggregation; after aggregation to the 2,655 NC census tracts, the resulting
matrix has at most $2655^2 \approx 7$M cells but is highly sparse.
The average NC home tract has commuting flows to approximately 15--30
distinct work tracts; most home-work tract pairs have zero flow.
This sparsity is exploited in the weighted Jaccard computation
(Section~\ref{sec:algorithm}).

\paragraph{Cross-state flows.}
LODES OD files are produced per state and record only intra-state commuting
flows.
Cross-state commuters (e.g., New Jersey residents commuting to New York City)
are partially captured in both states' files but with geographic coverage gaps.
For single-state redistricting runs, cross-state flow omission means that
border tracts may have artificially low commuting destination overlap with
tracts that would be captured if inter-state flows were included.
This is an acceptable limitation for the purpose of the M.4 edge weight:
the within-state commuting pattern still provides a meaningful community signal,
and border-state anomalies affect a small fraction of tracts.

\subsection{Functional Economic Areas and Commuting Sheds}

\paragraph{Tolbert-Sizer functional economic areas.}
\citet{tolbert1996labor} defined functional economic areas (FEAs) for the
entire United States by identifying county-level ``nodal clusters'' where
commuting flows indicated economic integration.
Their 1996 analysis produced 741 FEAs covering all U.S.\ counties.
The methodology has been updated periodically and provides the intellectual
foundation for Metropolitan Statistical Area delineation by the Office
of Management and Budget.

FEAs are defined at the county level; M.4 extends this concept to the
census tract level.
Two adjacent census tracts that share a commuting destination cluster belong to
the same functional economic community at the sub-county scale.
This captures a finer-grained community structure than FEA analysis:
within a single FEA, there may be multiple distinct sub-communities with
different primary employment destinations.

\paragraph{Research Triangle as a test case.}
The Research Triangle Park (RTP) area in North Carolina provides a well-studied
example of commuting shed community structure.
Workers residing in Raleigh, Durham, Chapel Hill, Cary, and Research Triangle
Park itself all commute to the same concentration of pharmaceutical,
biotechnology, and information technology employers located in the RTP
complex and the I-40/US-1 corridor.
This creates a natural commuting shed that spans four counties (Wake, Durham,
Orange, Chatham) but excludes the Charlotte metro tracts 100 miles to the
southwest, which commute to a completely different employment cluster
(finance in Uptown Charlotte, distribution in the I-85/I-77 corridor).

The M.4 algorithm is expected to assign high weighted Jaccard similarity to
Wake-Durham-Orange tract pairs that share RTP as a primary destination
and low similarity to pairs that span the RTP/Charlotte employment divide.
This prediction is the basis for the NC-14 polycentric separation test
in Section~\ref{sec:results}.

\paragraph{Milwaukee suburban commuting sheds.}
Waukesha County (west of Milwaukee) and the northern Milwaukee suburbs
(Ozaukee, Washington counties) both exhibit high commuting flows to
downtown Milwaukee and the western suburban employment corridor (Brookfield,
Wauwatosa, Menomonee Falls).
Despite not being geographically contiguous with each other, tracts in these
three counties share destination overlap through their common primary
destinations.
The M.4 algorithm is expected to assign similarity scores above 0.5 to
pairs of tracts from these counties that share Milwaukee-area primary
destinations.
